\begin{table}[H]
\centering
\caption{Linearity Test of the DiD Dose--Response: Connectivity and Firm Outcomes}
\label{tab:linearity_did}
\scriptsize
\begin{tabular}{llccccc}
\toprule\toprule
Outcome & Sample & $N$ & \shortstack{Bins\\(req.\ 50)} & \shortstack{Sup-$t$\\stat} & $p$-value \\
\midrule
Log Dec.\ wage & Intensive margin & $17{,}008$ & 50 & 3.015 & 0.243 \\
Log hourly wage & Intensive margin & $17{,}008$ & 50 & 2.261 & 0.852 \\
Log employment & Intensive margin & $17{,}120$ & 50 & 2.452 & 0.711 \\
\# CBA clauses & Intensive margin (CBA) & $10{,}273$ & 50 & 2.540 & 0.663 \\
\bottomrule
\end{tabular}
\begin{minipage}{\linewidth}
\footnotesize\vspace{4pt}
\textit{Notes:} This table reports Cattaneo, Crump, Farrell, and Feng (2024) sup-$t$ tests of the null hypothesis that the conditional expectation of each residualized outcome is a degree-1 polynomial (linear) in residualized treatment intensity $D_{jt}=Conn_j\times Post_t$. Both the outcome and $D_{jt}$ are residualized by Frisch--Waugh--Lovell on the same firm fixed effects, industry$\times$year, month-of-year$\times$year, and microregion$\times$year fixed effects, pre-treatment outcome, employment, and worker-flow quartile bins interacted with year, and the post indicator, as in the baseline intensive-margin specification. The sample is the balanced panel of untreated, positively connected firms (2009--2016; CBA periods for clause counts). Standard errors are clustered by firm; inference uses 2{,}000 simulation draws on a grid of 50 evaluation points per bin. $p$-values below $0.10$ are bolded.
\end{minipage}
\end{table}
